%!TEX program = lualatex
\documentclass[ngerman,headheight=70pt]{scrartcl}
\usepackage[ngerman]{babel}
\usepackage{blindtext}
\usepackage[utf8]{luainputenc}
\usepackage{scrlayer-scrpage}
\usepackage[T1]{fontenc}
\usepackage{fontspec}
\usepackage{multicol}
\usepackage{enumitem}
\usepackage[german=quotes]{csquotes}
\usepackage[bottom=2cm,footskip=8mm]{geometry}
\setmainfont{Source Sans Pro}
\MakeOuterQuote{"}

\setlength{\parindent}{0pt}
\setlength{\parskip}{6pt}

\pagestyle{plain.scrheadings}
\thispagestyle{scrheadings}
\lohead{\vspace*{-1.5cm}\includegraphics[scale=0.4]{../images/stupa_logo_rot.pdf}}
\rohead{\\Präsidium des\\Studierendenparlaments}

\lofoot[Universität Hamburg · Präsidium des Studierendenparlaments\\\hrule Von-Melle-Park 5 · D-20146 Hamburg]{Universität Hamburg · Präsidium des Studierendenparlaments\\\hrule Von-Melle-Park 5 · D-20146 Hamburg}
\cofoot[]{}
\cohead[\thepage]{}

\renewcommand{\thesubsection}{TOP \arabic{subsection}}
\renewcommand{\thesubsubsection}{\arabic{subsubsection}.}
\setcounter{subsection}{-1}

\begin{document}
    UHH · StuPa-Präsidium · Von-Melle-Park 5 · D-20146 Hamburg

    %%% PAGE ONE (EDIT LATER) %%%%
    % Hier den Titel des Protokolls eintragen
    \section*{Protokoll der 6. Sitzung des Studierendenparlaments vom 30. Juni 2016}

    % Protokollanten eintragen
    \textbf{Protokoll: Jim Martens}\\
    \textbf{Ort: S28 in der ehemaligen HWP}\\
    \textbf{Beginn: 18.15 Uhr}\\
    \textbf{Ende: Uhr}

    % Wohl die komplizierteste Stelle im Protokoll
    % Als Faustregel sollten nicht mehr als 3 Namen auf einer Zeile stehen
    % Einfach an dieser Vorlage orientieren
    % Im Idealfall ist hier nicht viel zu ändern, da kaum jemand fehlt.
    % Rücktritte müssen natürlich pro Protokoll angepasst werden
    \vspace{0.5cm}
    \begin{tabular}{ll}
        Anwesend: & \\
            RCDS (5 Sitze): & Ramon Weilinger, Antonia Niecke, Ramin Shakiba, \\
                            & Jennifer Maack \\
             CampusGrün (14 Sitze): & Laura Franzen, Geoffrey Youett, Philipp Droll, \\
                                   & Tahnee Herzig, Svenja Horn, Jim Martens,\\
                                   & Mirzo Ulugbek Khatamov, Armin Günther, Martin Sievert\\
                                   & Elena Rysikova, Yasemin Günther \\
             Bier-Liste (2 Sitze) : & Jakob Pape, Karen Martirosian \\
             WiWi (2 Sitze): & Claas-Friso Hente \\
             Unicorns (5 Sitze): & Johannes Peplow, Annkathrin Löffler, \\
                                 & Andreas Hartkamp, Marielle Hermstrüwer \\
             Liste LINKS (3 Sitze): & Gunhild Berdal, Till Petersen, Sinah Mielich \\
             HWP (2 Sitze): & Ajdina Karahasan \\
             MIN (4 Sitze): & \\
             SDS* (3 Sitze): & Mena Winkler, Jacob Petersein, Artur Brückmann \\
             Bart-LISTE (2 Sitze): & Timo Zeimet, Dominic Laumer \\
             LHG (1 Sitz): & \\
             harte zeiten (1 Sitz): & Tobias Berking \\
             Jura (1 Sitz): & \\
             AL (2 Sitze): & Karima Schulze, Henri Weber \\
            & \\
        Entschuldigt: & Benjamin Welling (RCDS),\\
                      & Melf Johannsen, Mario Moldenhauer (CG), Samet Gunay,\\
                      & Kay Zöllmer (WiWi), Katharina Kucza (Unicorns), Ailina Salten, \\
                      & Lotte Rullkötter, Thea Wahlers (alle MIN), Tobias Heisig (LHG)\\
                      &\\
        Unentschuldigt abwesend: & Jan Detampel (MIN), \\
                                & Johann Baumhoefener (Jura)\\
                                &\\
        Rücktritte: & Freya Schmitz (CG) \rightarrow Sandra Fietz\\
    \end{tabular}
    \newpage
    %%% Vorgeschlagene Tagesordnung %%%
    \underline{Vorgeschlagene Tagesordnung}
    \begin{enumerate}[label={\textbf{Top \theenumi}},leftmargin=*]
        % folgende Zeile bei Sitzungsfortsetzungen einkommentieren und x durch
        % die Zahl des zuletzt behandelten TOPs ersetzen
        %\setcounter{enumi}{x}
        \item Diskussion mit Wissenschaftssenatorin Katharina Fegebank (60 Min.)
        \item Haushalt 2016/2017 (V1617-025, V1617-030) (75 Min.)
        \begin{enumerate}
            \item 2. Lesung (70 Min.)
            \item 3. Lesung (5 Min.)
        \end{enumerate}
        \item Nachtragshaushalt 2015/2016 (V1617-026, V1617-030) (25 Min.)
        \begin{enumerate}
            \item 2. Lesung (20 Min.)
            \item 3. Lesung (5 Min.)
        \end{enumerate}
        \item Feministischer Juristinnentag 2017 in Hamburg (V1617-028) (30 Min.)
        \item Weiter initiativ mit der Strategie "Uni für Alle statt Exzellenz" (V1617-029) (30 Min.)
        \item Verlegung der StuPa-Wahl in das Sommersemester (V1617-015) (60 Min.)
        \item Verschiedenes
    \end{enumerate}

    \newpage

    % Hier beginnt Teil A. Bei Fortsetzungen kommt hier nur der Teil über
    % die Feststellung der endgültigen Fassung der Tagesordnung.
    % Für sonstige Sitzungen ist dies bereits das Gerüst für die Formalia.
    % Man muss nur noch die protokollierten Inhalte eintragen.

    {\Large \underline{Teil A}}

    \subsection{Formalia}

    \textit{Gunhild eröffnet die Sitzung um 18.24 Uhr.}

    \subsubsection{Geschäftsbericht Präsidium}

    Haben uns Gedanken über Diskussion gemacht. Exzellenzstrategie und
    Finanzierung der Hochschulen werden in ein Themenblock zusammengefasst.
    Senatorin wird kurz etwas vortragen, das Ziel ist schnell in Diskussion zu
    kommen.

    Wir haben vor uns in den nächsten Wochen Gedanken über die Wahlen
    zu machen. Auch im Bereich Öffentlichkeitsarbeit steht noch einiges an.

    \subsubsection{Anfragen an das Präsidium}

    Keine Fragen an das Präsidium.

    \subsubsection{Geschäftsbericht AStA}

    Franzi hält den Geschäftsbericht des AStA. Letzte Woche Do war
    Wissenschaftsausschusssitzung. Artur und Moritz waren auf der Senatsbank
    und konnten einige gute Punkte fürs Studierendenwerk anbringen.

    Das Em-Begleitprogramm kommt gut voran. In dem Zusammenhang Treffen mit
    franz. Studis.

    6.7.
    9.7. Here to stay
    14.7. em und nationalismus

    AstA hat Exzellenzstrategie ausgewertet.
    Am 7.7. gibts Veranstaltung zu Ex-Ini mit libertären Mauerkaninchen.
    Alternativen zu Drittmittel.

    Studentische Arbeit gegen Rechts. Beteiligung an Anti-Afd-Veranstaltung
    in Rahlstedt. Flugblätter wurden verteilt.
    Am 5.7. findet Abend gg Rechts statt, veranstaltet von Ausschuss gg rechts.

    AStA-Klausur zur Frage "Wie gegen rechts arbeiten" wird in den nächsten
    Wochen stattfinden.

    18:30

    Golnar berichtet aus dem Ris. Russische Linke und ihre Lage im aktuellen
    Russland von Putin. Ein Metallarbeiter war zu Gast. Veranstaltung war
    im HWP-Hörsaal und gut besucht. Es wird politische VV stattfinden. Akute
    Abschiebungsbedrohung von vielen ausländischen Studierenden.
    Nächsten Monat gibt es Termin mit Harburger Ausländerbehörde.

    Meta: Geoffrey steckt das VGA-Kabel um. Geoffrey führt ein Video von der
    Anti-ExIni-Demo vor.

    \subsubsection{Anfragen an den AStA}

    Till fragt, ob sich die Studierenden der franz. Gewerkschaft zum Verhältnis
    von gewerkschaft zu Studibewegung geäußert haben und inwieweit das für das
    Rotstiftbündnis angewendet werden kann.

    Artur: In Frankreich gibt es nicht an jeder Hochschule eine pol. Gruppe, sondern
    nationale gewerkschaften. Größte ist unef. Sie begreifen Studierende als
    intellektuelle Arbeiter. Pflegen Nähe zu Gewerkschaften. Es wurde sich
    intensiv mit Arbeitsgesetz auseinandergesetzt. Nach langer Phase von Beschäftigung
    mit Gesetz wurde breites Bündnis gegen das Gesetz gegründet. Studi hat gesagt,
    dass bei bereits abflauenden Aktivitäten alle Gewerkschaften 2 Demos pro
    Woche machen sollen.

    Golnar: Aktionsformen der Studis in Frankreich unterscheiden sich nicht
    grundlegend von unseren Aktivitäten. Es ist somit Haltungsfrage. Sie machen
    sich nicht Gedanken, ob es die konkrete Klientel etwas angeht. Es sei
    Gesellschaftspolitik und geht somit alle was an. Sie würden sich zu wenig
    mit Militarismus auseinandersetzen und seien zu sehr aktionistisch unterwegs.
    Positivorientierung fehlt.

    Till fragt: Gab es besondere Erkenntnisse bei den EM-Alternativ-Veranstaltungen?

    \subsubsection{Dringlichkeitsanträge des AStA}

    Keine.

    \subsubsection{Aktuelle Stunde (falls entsprechender Antrag vorliegt)}

    Keine.

    \subsubsection{Feststellung der endgültigen Fassung des Teils B der Tagesordnung}

    Vorziehen von Nachtragshaushalt vor Haushalt. Kein Widerspruch.

    %%% BESCHLOSSENE TAGESORDNUNG %%%
    \underline{Endgültige Fassung der Tagesordnung}
    \begin{enumerate}[label={\textbf{Top \theenumi}},leftmargin=*]
        % folgende Zeile bei Sitzungsfortsetzungen einkommentieren und x durch
        % die Zahl des zuletzt behandelten TOPs ersetzen
        %\setcounter{enumi}{x}
        \item Diskussion mit Wissenschaftssenatorin Katharina Fegebank (60 Min.)
        \item Nachtragshaushalt 2015/2016 (V1617-026, V1617-030) (25 Min.)
        \begin{enumerate}
            \item 2. Lesung (20 Min.)
            \item 3. Lesung (5 Min.)
        \end{enumerate}
        \item Haushalt 2016/2017 (V1617-025, V1617-030) (75 Min.)
        \begin{enumerate}
            \item 2. Lesung (70 Min.)
            \item 3. Lesung (5 Min.)
        \end{enumerate}
        \item Weiter initiativ mit der Strategie "Uni für Alle statt Exzellenz" (V1617-029) (30 Min.)
        \item Verlegung der StuPa-Wahl in das Sommersemester (V1617-015) (60 Min.)
        \item Verschiedenes
    \end{enumerate}

    TO wurde angenommen.

    \subsubsection{Feststellung der Beschlussfähigkeit}

    28. ParlamentarierInnen anwesend und beschlussfähig.

    \subsubsection{Genehmigung der Protokolle der vorangegangenen Sitzungen}

    Protokoll ist genehmigt.

    19:00

    \vspace{0.5cm}
    {\Large \underline{Teil B}}

    % Teil B beginnt
    % Dieser wird sich für jede Sitzung unterscheiden.
    % Jeder TOP ist eine subsection.
    % Vorhandene Beispiele anfangs entfernen und durch tatsächliche TOPs ersetzen

    % Folgende Zeile bei Sitzungsfortsetzungen einkommentieren und x mit
    % der Zahl des zuletzt behandelten TOPs ersetzen.
    %\setcounter{subsection}{x}
    \subsection{Diskussion mit Wissenschaftssenatorin Katharina Fegebank}

    Fegebank trifft 3 min vor 19 Uhr ein.

    Fegebank begrüßt und macht ein paar einleitende Worte. Sie hat zur Kenntnis
    genommen, dass vom AStA und StuPa Unterstützung für die Hamburger Haltung
    gegen ExIni gegeben wurde. Sie klärt Missverständnis zu ihrer Position auf.
    Sie war nie gegen ExIni, sondern nur gegen konkrete Ausgestaltung. Fegebank
    feiert die Verhandlung bzgl. Exini als Erfolg für gesamte Hochschullandschaft
    und Standort Hamburg. Hamburger Uni könne mit viel Selbstbewusstsein in die
    nächste Runde gehen. Uni habe schon 2 Exzellenzcluster.

    Kurzum: Exzellenzwerbung und warum das alles geil ist. Neoliberales Blabla.
    Keine Substanz, sondern nur Redephrasen.

    Dies sei Einstieg in die Erhöhung der Grundfinanzierung und würde häufig
    über 3\%, manchmal 4\% entsprechen. Sie geht auf destruktive Wirkung der
    Exini ein. Forschungsförderung soll diese Binnenwirkung der Exini verhindern.

    Es werden Gespräche für "Aufwuchs" in TUHH und HAW geführt. HSP-Mittel sollen
    verstetigt werden.

    --

    Gunhild dankt für einführende Worte.

    --

    Till führt an, dass es Kritik aus Uni gab entgegen Aussage von HR.
    Er spricht FAZ-Artikel an.
    Warum muss Exzellenz sein und schafft nicht für alle eine bessere
    Bedingung? Welches Verständnis von wissenschaftlicher Qualität steckt
    dahinter? Woher kommt Verständnis, dass Wissenschaft am besten funktioniert,
    wenn es Gewinner und Verlierer gibt?

    19:15

    Jochen erzählt über Berliner TU. In Hamburg ist TU sehr stark
    ingenieurwissenschaftlich. Wenn TU wächst, sollte die gesellschaftliche
    Bedeutung von Technologie evaluiert werden können.

    Franziska weist auf den Double Speak hin. Antragsprosa aber hinter
    vorgehaltener Hand wird sich stark beschwert. Im Imbodenbericht wird
    erwähnt, dass Wimis in ihren produktivsten Jahren sich mit Antragsschreiben
    und hinterherhecheln von wenigen Festanstellungen beschäftigen.

    Laura fragt, wie man sich an einer Uni mit Exini und wenigen Festanstellungen
    eine Familie vorstellen kann.

    Fegebank findet es nicht schlecht, dass es Wettbewerb gibt. Sehr viel
    Standortideologie. Dissens bleibt wohl bestehen.

    19:30

    Leistung, dass Elbe in HH bereits im Juni über die Ufer tritt.
    Sie arbeitet an Entfristung.

    Lehre als Alleinstellungsmerkmal?

    --

    Tobias stellt die Notwendigkeit von Exini in Frage. Bewerbungsverfahren
    ist Erziehungsverfahren. Wissenschaft muss Bestehendes in Frage stellen.
    Warum Wettbewerb und nicht Kooperation?

    Gunhild erinnert sich an Brecht. Besser im Wettbewerb ist nicht unbedingt
    besser für Menschheit. Gunhild versteht Zusammenhang zwischen Vielfalt
    und Exzellenz nicht. Drittmittel sind ein Problem bei Unabhängigkeit von
    Forschung.

    Artur weist darauf hin, dass die Exini eine Mainstreamisierung verursacht.
    Hauptsächlich wird auf Mittel geachtet, die Inhalte sind nachgelagert.

    Golnar merkt an, dass es ein Demokratiedefizit in dem Verfahren gibt.
    Wie steht Fegebank zu Demokratiedefizit? Golnar fragt, wie Vielfalt und
    Exini zusammenpassen.

    19:46

    Golnar weist auf Ambivalenz im deutschen Hochschulsystem hin.

    Ramon befürwortet Exini. Ist Exzellenzbegriff nicht zu nat.wiss.lastig.
    War Kritik nicht zu spät?

    Fegebank sagt, dass HH seit Beginn an Kritik geäußert hat. Sie hat noch keine
    Infos zu Wissenschaftsratsbegehung.

    Sie entgegnet kritische Anmerkungen mit Frage, was denn passieren würde, wenn
    man der Uni auf einen Schlag 0,5 Mrd. geben würde.
    Sichtbarkeit spielt eine große Rolle bei der Hochschulfinanzierung.
    Es gibt Leute, die hinterfragen, warum so viel Geld in die Hochschulen
    geht.

    Ailina fragt über Arbeitssicherheit und Kosten für Investitionen.
    Es gibt in der Biologie viele kaputte Labore.

    Till zieht Zwischenfazit.

    Fegebank schaut sich gerne die Biologie an.

    20:08

    Pause bis 20:20, verlängert bis 20:25

    \subsection{Nachtragshaushalt 2015/2016}
    \subsubsection{2. Lesung}

    Ramon erklärt das Prozedere für die Abstimmungen. Melf stellt das erste
    Paket für den NHH vor.

    Keine Gegenrede zum Paket.

    Einstimmig wird der Beschlussempfehlung gefolgt.

    Melf stellt das zweite Paket vor. Keine Gegenrede. Einstimmig wird der
    Beschlussempfehlung gefolgt.

    \subsubsection{3. Lesung}

    M:1:0 wird NHH angenommen.

    \subsection{Haushalt 2016/2017}

    \subsubsection{2. Lesung}

    Melf stellt die Pakete für den Haushalt 16/17 vor. Melf stellt das erste
    Paket für den Haushalt genauer vor. Nico hält die Fürrede für den Antrag
    den Topf für PKs zu erhöhen.

    Tobias hält Fürrede für Antrag den Topf für ReferentInnen zu erhöhen. Er
    steigt ein in Debatte.

    20:40

    Karima fragt Till, ob er sich nicht um die Pluralität Sorgen macht, wenn
    7 1/2 PK-Stellen für 2 Ref-Stellen wegfallen.

    Till möchte nicht, dass bei PKs und Refs über Stellen gesprochen wird.
    Er ist der Meinung, dass Geld auf Dauer keine Lösung ist.

    Tobias verweist auf AStA in Lüneburg, der keine PKs hat.

    M:5:4 für Beschlussempfehlung.

    Melf stellt zweites Paket vor. Er stellt seine Meinung zum ÄA zu Aktionen
    vor. Er stellt auch den Punkt Vereinsmitgliedschaften FSRe vor.

    Tobias hält Fürrede für Antrag zu Aktionen und Veranstaltungen der VS.

    Melf hält Plädoyer für einen realistischen Haushalt, wo nur so viel
    Geld in Posten gesteckt wird, wie benötigt wird.

    Till hält Plädoyer für eine Erhöhung des Aktivitätentopfes.

    M:8:2 für Beschlussempfehlung.

    Melf stellt das dritte Paket FZS vor und hält eine Fürrede. Vincent hält
    eine Gegenrede zu den Anträgen und sagt, dass Vollmitglied = Vollmitglied
    bedeuten muss. Erhöhung im Personalbereich aber sinnvoll.

    Geoffrey pflichtet Vincent bei, dass es sinnvoll wäre einen Ref zu haben.
    Der Betrag sollte nicht verringert werden. Er geht auch auf das Argument ein,
    dass andere Studierendenschaften weniger zahlen. Dies lässt sich nicht
    eindeutig verifizieren.

    21:00

    Artur geht auf den europäischen Zusammenhang ein.

    Till führt eine Schauspielleistung vor und sagt, dass Geschacher reines
    Marktprinzip und daher abzulehnen ist.

    Nach Protesten gegen Rechte in Sachsen wurden die Studierendenschaften
    in Sachsen mit den Stimmen der NPD verboten und degradiert auf reine
    Vereinsebene.

    6:M:2 gegen Beschlussempfehlung.
    7:M:2 gegen Antrag Mitgliedsbeitrag zu senken.
    1:14:viele gegen Erhöhung Aufwandsentschädigung.

    Melf stellt viertes Paket vor. Till bittet in Zukunft zu bedenken Begründungen
    so zu verfassen, dass sie verständlich sind.

    M:0:1 für Beschlussempfehlung.
    M:1:4 für Haushalt

    21:17

    Pause bis 21:30, 21:35

    \subsection{Weiter initiativ mit der Strategie "Uni für Alle statt Exzellenz"}

    \textit{Fraktionspause Liste Links 3 Minuten}

    Till stellt den Antrag vor.

    21:45

    Artur meint, dass es eine weitere Entstellung zur Kenntlichkeit
    gegeben hat.
    Differenz zwischen "Trockenschwimmen" und sinnvoller Forschung muss
    herausgestellt werden. Fegebanks Punkt über die Sichtbarkeit der Uni hat
    einen wahren Kern. Allerdings sollte es in der Sichtbarkeit nicht um
    Sonderforschungsbereiche, sondern um qualitativ sinnvolle Aktionen wie z.B.
    Engagement für Flüchtlinge gehen.
    Man muss besser klarmachen, dass eine Alternative zur Exini real
    möglich ist. Man könnte von der Bewegung um Bernie Sanders lernen.

    Tobias stellt historische Zusammenhänge dar.

    Armin meint, dass man generell das Drittmittelsystem bekämpfen müsste.

    Till befürwortet eine solche allgemeine Kritik. Er geht auf die Diskussion
    mit Fegebank ein. Sie soll deutlich ins Schwimmen geraten sein. Vor allem
    habe sie mit dem Versprechen von 500 Mio. Euro einen Fehler begangen, den
    man als Senatorin nicht machen sollte. Abfolge:
    An sich ist Mittelerhöhung richtig, oh das darf ich nicht sagen, habe ich
    gar nicht erst gemeint.
    Till würde gerne die Diskussion mit Fegebank auswerten.

    Mario unterstützt Ansicht von Armin und möchte vor Abstimmung genauen
    Text wissen.

    Geoffrey findet ihr Beispiel mit den 500 Mio. Euro zynisch, dass das Geld
    zu wenig wäre und die aktuelle Situation daher besser. Vor allem die Profilbildung
    macht ihm Sorgen.

    22:05

    Yasemin hat den Eindruck, dass sie sehr unter Druck stand. Sie ist ein
    wenig enttäuscht von dem Besuch. Fegebank hat überhaupt nicht
    gesagt, dass soziale Wissenschaften wichtig sind.

    Jim hält einen Redebeitrag. Fegebank hat mit neoliberalem Werbebeitrag
    begonnen und ist langsam in eine echte Diskussion eingestiegen.
    Sie schien ein bisschen gefangen zu sein zwischen Senatsvertretung und
    persönlicher Ambition.

    Laura äußert sich über die Profiloberstufe und wie desaströs das gelaufen
    ist. Fegebank hat kein einziges Mal Effizienz gesagt, obwohl sie eigentlich
    nichts anderes getan hat.

    Armin: Fegebank hat häufig betont, dass in dem Haushalt eine Erhöhung
    der Grundfinanzierung um 3\% bis 4\% geplant sei, obwohl die Summe lächerlich
    ist. Man sollte sie zu einer Veranstaltung zum Clisep? einladen.

    Golnar weist darauf hin, dass sie am Anfang auf Till entgegnete, dass
    man eben zwei unterschiedliche Religionen habe. Sie hat dann aber herausgefunden,
    dass an der Uni mit echten Argumenten gearbeitet wird.
    Man dürfe jetzt nicht mit dem Kämpfen aufhören, obwohl es kleine Geldgeschenke
    gibt.

    Till verweist auf die Demo gegen Exini und den Zusammenhang zwischen
    Rathaus und Börse/Handelskammer. Fegebank sei eine junge Senatorin und könne
    sich dem Argument nicht komplett verschließen. Außerdem sei sie sehr
    druckanfällig für beide Seiten.

    Till stellt ÄA vor.

    Artur stellt seinen ÄA vor.

    22:30

    Till erklärt historischen Werdegang von Rap und HipHop.

    M:2:0 für Annahme des Antrags.

    \subsection{Verlegung der StuPa-Wahl in das Sommersemester}

    \textit{wurde zurückgezogen}

    \subsection{Verschiedenes}

    Ramon fragt, ob es den Wunsch gibt am 14. Juli eine Stupa-Sitzung zu machen.

    Franzi spricht sich für eine Sitzung des Parlaments aus.

    Melf sieht keine Notwendigkeit für eine Sitzung.

    Till spricht sich für eine Sitzung aus. Er sieht keinen Zusammenhang
    zwischen Franzis Aussage und Melfs Aussage. Man sollte eine Sitzung machen,
    um noch einmal viel vor dem WS zu machen.

    Jim spricht sich grundsätzlich für eine Sitzung aus, sieht aber Probleme
    bei der Beschlussfähigkeit.

    Tobias meint, dass eine solche Sitzung mit Beschlussfähigkeit möglich ist,
    wenn man es darauf anlegt.

    Gunhild spricht sich für eine Sitzung aus und findet, dass Personen, die nicht
    kommen können, sich nicht gegen eine Sitzung aussprechen sollten.
    Gunhild appelliert, dass die Ausschüsse wieder mehr arbeiten sollten.

    Melf spricht sich erneut gegen eine Sitzung aus, findet aber einzelne
    Treffen in der VL-freien Zeit sinnvoll.

    Till findet es wichtig, dass man es auf eine Sitzung anlegt.

    Ramon schlägt vor ein Stimmungsbild zu machen. 10:12:1 für eine weitere Sitzung.

    Ramon schließt um 22:57 die Sitzung.

\end{document}
